\section{Ingeniería de visualización y monitoreo}

\subsection{Diseño del dashboard}
Cada módulo de Compound AI debe tener un indicador visible en el panel de operaciones: la trace fidelity curve se compara horizontalmente, la judge reward curve y la user satisfaction signal se resumen en una vista de dos en uno, y el solver output resalta en otro panel la constraint violation rate. La siguiente tabla enumera los dashboards clave y la frecuencia de actualización recomendada.
\begin{table}[H]
\centering
\begin{tabular}{lp{8cm}r}
\toprule
Dashboard & Métricas incluidas & Frecuencia de actualización recomendada \\
\midrule
Trace health & fidelity, cost variance, skew coverage & 1 minuto \\
Judge stability & distribución de accuracy/proactivity/credibility, NPS proxy scores & 30 segundos \\
Constraint guard & solver violation rate, brush 工艺 compliance & 10 segundos \\
\bottomrule
\end{tabular}
\caption{Principales dashboards de Compound AI y estrategia de actualización}
\end{table}
\begin{knowledgebox}{El ritmo de retroalimentación de la visualización}
Escalona la frecuencia de actualización de los distintos módulos de monitoreo según su criticality, para evitar que actualizar todas las métricas al mismo tiempo sature el monitoreo; al mismo tiempo, usa \texttt{trace\_id} como dimensión común para enlazar en el dashboard la información de trace, judge y solver, lo que ayuda a los ingenieros a investigar entre módulos.
\end{knowledgebox}

\subsection{Repetición de trace y plan}
La verdadera transparencia proviene de la capacidad de repetición de trace/plan: después de ejecutar cada plan se escribe un \texttt{trace\_id.json} que contiene la puntuación heuristic, las violaciones de restricciones y el judge reward, de modo que en la página web de operaciones se puede reproducir paso a paso el search trace, el judge dialogue y el solver output. Para que cada repetición tenga una impresión visual, escribimos las marcas de tiempo de los frames clave en el mismo log, de manera que al reproducir se puedan cargar la captura de la slide o la imagen del frame correspondientes.
\begin{knowledgebox}{La repetición acelera la auditoría}
Una repetición de trace no solo permite reconstruir el decision path, sino que además le permite al engineer ver directamente el visual evidence alineado con la slide. Por ejemplo, la slide de APAC en el segmento 00:16:08--00:16:22 aparece en la repetición como anotación del judge trigger; el diagrama de arquitectura de DSPy en el segmento 00:32:25--00:32:45 se usa para explicar el cálculo paso a paso del solver.
\end{knowledgebox}
\begin{warningbox}{Los datos de repetición no pueden faltar}
La falta de alineación temporal entre trace y frame retrasa la investigación. Es indispensable escribir en el log al mismo tiempo \texttt{trace\_id}, \texttt{frame\_path} y \texttt{timestamp}; de lo contrario, al reproducir solo se podrá comparar de manera aproximada y el juicio operativo se distorsionará.
\end{warningbox}

\subsection{Resumen de la sección}
El punto central del monitoreo de ingeniería es fusionar las metrics, el visual evidence y la lógica de repetición en tres cadenas consultables como una sola unidad, de modo que cuando ocurra un evento se pueda localizar rápidamente el trace frame, el judge dialog y el solver plan correspondientes, y así responder con precisión a las fallas.
